% Clase 4 — Razonamiento bajo Incertidumbre (parte 2)
% Lectura: AIMA Cap. 12
\section{Razonamiento bajo Incertidumbre II}

\subsection{Inferencia por Enumeración}

Para cualquier consulta $P(X \mid \mathbf{e})$:

$$P(X \mid \mathbf{e}) = \alpha \sum_{\mathbf{y}} P(X, \mathbf{e}, \mathbf{y})$$

donde $\mathbf{y}$ son las variables ocultas y $\alpha$ es el normalizador.

\subsection{Variables y Modelos}

\begin{itemize}
    \item \textbf{Variable aleatoria}: puede ser discreta o continua.
    \item \textbf{Distribución de probabilidad}: asigna probabilidades a todos los valores.
    \item \textbf{Tabla de probabilidad condicional} (CPT): especifica $P(X \mid \text{Parents}(X))$.
\end{itemize}

\begin{example}[Enumeración sobre Dolor de muela, Caries y Detector]
    Sea la distribución conjunta completa $P(\text{Caries}, \text{Dolor}, \text{Detector})$:
    \begin{table}[H]
        \centering
        \begin{tabular}{lcccc}
            \toprule
             & \multicolumn{2}{c}{Dolor} & \multicolumn{2}{c}{$\neg$Dolor} \\
             & Detector & $\neg$Detector & Detector & $\neg$Detector \\
            \midrule
            Caries      & 0.108 & 0.012 & 0.072 & 0.008 \\
            $\neg$Caries & 0.016 & 0.064 & 0.144 & 0.576 \\
            \bottomrule
        \end{tabular}
        \caption{Distribución conjunta completa (suma de las 8 celdas = 1).}
    \end{table}
    Para calcular $P(\text{Caries} \mid \text{Dolor})$ por enumeración, se \textbf{suma sobre la variable oculta}
    Detector, dejando fija la evidencia Dolor:
    \[
        P(\text{Caries} \mid \text{Dolor}) = \alpha \sum_{\text{det}} P(\text{Caries}, \text{Dolor}, \text{det})
        = \alpha \, (0.108 + 0.012) = \alpha \, (0.12)
    \]
    De forma análoga, $P(\neg\text{Caries} \mid \text{Dolor}) = \alpha\,(0.016+0.064) = \alpha\,(0.08)$. Normalizando
    con $\alpha = 1/(0.12+0.08) = 5$: $P(\text{Caries}\mid\text{Dolor}) = 0.6$ y $P(\neg\text{Caries}\mid\text{Dolor}) =
    0.4$. Este es el patrón general de inferencia por enumeración: fijar la evidencia, sumar sobre las variables
    ocultas, y normalizar al final.
\end{example}
